\chapter{Methodology}
\label{chap:methodology}

\section{Data Sources}

The empirical workflow combines two public data sources: Freddie Mac sample mortgage loan files and monthly macroeconomic series from the Federal Reserve Economic Data (FRED) database \citep{freddiemac_loanlevel,fred_database}. The mortgage files provide loan-level origination characteristics and monthly servicing performance records. The macroeconomic data provide the aggregate state variables used to infer the latent regime process.

\begin{table}[htbp]
    \centering
    \caption{Data sources and empirical role}
    \label{tab:data-sources}
    \begin{tabular}{L{0.24\textwidth}L{0.32\textwidth}L{0.29\textwidth}}
        \toprule
        Source & Main variables & Role in the methodology \\
        \midrule
        Freddie Mac sample origination files & Borrower and loan characteristics at origination, including credit quality, leverage, product type, and loan terms. & Define the cross-sectional risk profile of each mortgage. \\
        Freddie Mac sample servicing files & Monthly performance status, delinquency indicators, remaining balance, current activity, and loan age. & Construct the loan-month panel and forward-looking default label. \\
        FRED macroeconomic series & Inflation, policy rates, unemployment, house prices, mortgage rates, Treasury rates, claims, payrolls, real income, financial conditions, and sentiment. & Construct interpretable macro factors, estimate the hidden regime, and produce filtered stress probabilities. \\
        \bottomrule
    \end{tabular}
\end{table}

Let $i=1,\ldots,N$ index loans and let $t=1,\ldots,T$ index calendar months. The raw mortgage records are transformed into a panel in which each observation corresponds to a loan-month. The origination file supplies time-invariant or slowly varying underwriting characteristics, while the servicing file supplies monthly performance information. The macroeconomic variables are merged by calendar month so that all loans observed in month $t$ receive the same macro vector.

The raw transformed macro vector is
\begin{equation}
    M_t=(M_{t,1},\ldots,M_{t,q})^{\top}\in\mathbb{R}^{q},
    \qquad t=1,\ldots,T.
    \label{eq:macro-vector}
\end{equation}
Its components are economically meaningful levels, changes, spreads, or growth rates. The candidate set comprises CPI inflation, the Federal Funds rate, unemployment and its change, house-price growth, the 30-year mortgage rate, the mortgage--Treasury spread, the 10-year Treasury rate and its change, the 10-year minus 2-year yield spread, initial claims and their change, payroll growth, real disposable-income growth, the National Financial Conditions Index, consumer sentiment, and sentiment change. Weekly series are aggregated to monthly frequency before the merge.

To avoid fitting a high-dimensional full-covariance HMM directly, the variables are oriented so that larger values represent greater stress and grouped into four transparent factors: housing/collateral, rate/refinancing, labour/income, and financial conditions. If $\mathcal{G}_g$ is the set of source variables assigned to group $g$ and $a_j\in\{-1,1\}$ is the adverse-direction sign, then
\begin{equation}
    F_{t,g}
    =
    \frac{1}{|\mathcal{G}_g|}
    \sum_{j\in\mathcal{G}_g}
    \frac{a_jM_{t,j}-\overline{a_jM_j}}{s(a_jM_j)},
    \qquad g=1,\ldots,r.
    \label{eq:interpretable-factor}
\end{equation}
Equal weighting preserves a clear mapping from source indicators to HMM inputs and prevents a group with many correlated series from dominating by construction. The factor vector $F_t=(F_{t,1},\ldots,F_{t,r})^{\top}$, with $r=4$ when all groups are available, is standardised again before HMM estimation:
\begin{equation}
    Z_{t,g}
    =
    \frac{F_{t,g}-\bar F_g}{s_g},
    \qquad g=1,\ldots,r.
    \label{eq:macro-standardisation}
\end{equation}
The current expanded empirical pass estimates these quantities and the HMM parameters over January 2014--September 2024, then computes the state recursion forward through that window. The reported state probabilities are therefore filtered with respect to observations but conditional on full-window transformation and HMM parameters. A strict real-time robustness design must re-estimate the factors, scaler, and HMM at each rolling origin; the manuscript retains this distinction explicitly rather than claiming that the present HMM is fully out of sample.

\section{Unit Of Analysis}

The unit of observation is one mortgage at one month. For loan $i$ in month $t$, the full observation is written as
\begin{equation}
    \mathcal{O}_{i,t}
    =
    \left(Y_{i,t}^{(12)}, X_{i,t}, M_t, \hat p_t^{\mathrm{Stress}}\right),
    \label{eq:loan-month-observation}
\end{equation}
where $Y_{i,t}^{(12)}$ is the 12-month serious-delinquency label, $X_{i,t}$ is the borrower and loan feature vector, $M_t$ is the observed macro vector, and $\hat p_t^{\mathrm{Stress}}$ is the filtered probability that month $t$ is in the stress regime.

The borrower and loan feature vector is represented as
\begin{equation}
    X_{i,t}
    =
    \left(X_{i,t,1}, X_{i,t,2}, \ldots, X_{i,t,p}\right)^\top
    \in \mathbb{R}^{p}.
    \label{eq:feature-vector}
\end{equation}
The intercept is modelled separately in the logistic equations. The components of $X_{i,t}$ consist of origination variables and current performance variables observed at month $t$. The feature list is fixed before model estimation and excludes any variable measured after month $t$. The construction follows the information principle that only information available at or before month $t$ may be used to predict the event over months $t+1$ through $t+12$.

The empirical panel is defined over the at-risk population. A loan-month is at risk when the loan is active and not already seriously delinquent at month $t$. Let $D_{i,t}$ denote delinquency at month $t$, $T_i^{\max}$ the last observed month for loan $i$, and $T_i^{\mathrm{term}}$ its first terminal-servicing month, if any. The implemented follow-up indicator is
\begin{equation}
    W_{i,t}^{(12)}
    =
    \mathbf{1}\{t+12\leq T_i^{\max}\}
    \mathbf{1}\{T_i^{\mathrm{term}}\text{ is missing or }T_i^{\mathrm{term}}>t+12\}.
    \label{eq:forward-window-indicator}
\end{equation}
Let $\mathcal{A}_{i,t}$ denote the event that the loan is active and eligible for forward prediction. The modelling sample is restricted to
\begin{equation}
    \mathcal{R}
    =
    \left\{(i,t): \mathcal{A}_{i,t}=1,\; D_{i,t}<90,\; W_{i,t}^{(12)}=1\right\}.
    \label{eq:risk-set}
\end{equation}
This restriction requires the observed servicing history to extend beyond the twelve-month horizon and excludes loans terminating inside that horizon. It prevents the model from learning from observations already in the target state and avoids treating mechanically truncated follow-up as a non-event.

\section{Primary Label}

The primary response variable is a 12-month transition to serious delinquency, defined as reaching 90 or more days past due within the next 12 months. For each at-risk loan-month $(i,t)\in\mathcal{R}$, the label is
\begin{equation}
    Y_{i,t}^{(12)}
    =
    \mathbf{1}
    \left\{
        \max_{h=1,\ldots,12} D_{i,t+h} \geq 90
    \right\}.
    \label{eq:primary-label}
\end{equation}
This definition aligns the statistical target with a forward-looking PD concept: the question is not whether the loan is currently distressed, but whether it transitions into serious delinquency over a fixed prediction horizon.

The conditional probability of interest is
\begin{equation}
    \pi_{i,t}
    =
    \Pr\left(Y_{i,t}^{(12)}=1
    \mid \mathcal{F}_t\right),
    \label{eq:conditional-pd}
\end{equation}
where $\mathcal{F}_t$ denotes the information available at month $t$. In the borrower-only baseline, $\mathcal{F}_t$ contains loan-level information only. In the macro-augmented baseline, it also contains observed macro variables. In the regime-aware model, it contains the filtered stress probability inferred from the macro HMM.

The 90+ days-past-due threshold is used because serious delinquency is a standard operational default proxy in credit risk and is closely related to regulatory default definitions based on material arrears \citep{bcbs_qis3_default,fed_basel2_default}. The fixed 12-month horizon creates a consistent target across the panel and supports direct comparison of predicted probabilities with realised default frequencies.

\section{Macro Regime Inference}

The latent macro regime is modelled using a two-state Gaussian hidden Markov model. Let
\begin{equation}
    R_t \in \{1,2\}
    \label{eq:regime-state-space}
\end{equation}
denote the unobserved macroeconomic regime in month $t$. After estimation, the two states are interpreted economically and labelled as Expansion and Stress according to their fitted macro profiles.

The regime process is assumed to follow a first-order homogeneous Markov chain:
\begin{equation}
    \Pr(R_t=k \mid R_{t-1}=j, R_{t-2},\ldots,Z_{1:t-1})
    =
    \Pr(R_t=k \mid R_{t-1}=j)
    =
    p_{jk},
    \label{eq:hmm-transition}
\end{equation}
for $j,k\in\{1,2\}$. The transition matrix is
\begin{equation}
    P
    =
    \begin{bmatrix}
        p_{11} & p_{12} \\
        p_{21} & p_{22}
    \end{bmatrix},
    \qquad
    p_{jk}\geq 0,
    \qquad
    \sum_{k=1}^{2}p_{jk}=1.
    \label{eq:transition-matrix}
\end{equation}
The diagonal entries $p_{11}$ and $p_{22}$ measure regime persistence. The off-diagonal entries measure transition probabilities between expansion and stress.

Conditional on the latent regime, the standardised macro vector follows a multivariate Gaussian distribution:
\begin{equation}
    Z_t \mid R_t=k
    \sim
    \mathcal{N}(\mu_k,\Sigma_k),
    \qquad
    k=1,2.
    \label{eq:gaussian-emission}
\end{equation}
The complete HMM parameter vector is
\begin{equation}
    \Theta
    =
    \left(\delta, P, \mu_1,\mu_2,\Sigma_1,\Sigma_2\right),
    \label{eq:hmm-parameter-vector}
\end{equation}
where $\delta=(\delta_1,\delta_2)$ is the initial regime distribution.

For an observed factor path $Z_{1:T}$ and a latent regime path $R_{1:T}$, the joint density is
\begin{equation}
    f(Z_{1:T},R_{1:T};\Theta)
    =
    \delta_{R_1}
    \left[
        \prod_{t=2}^{T}p_{R_{t-1},R_t}
    \right]
    \left[
        \prod_{t=1}^{T}
        \phi_r(Z_t;\mu_{R_t},\Sigma_{R_t})
    \right],
    \label{eq:hmm-joint-density}
\end{equation}
where $\phi_r(z;\mu,\Sigma)$ denotes the $r$-dimensional Gaussian density:
\begin{equation}
    \phi_r(z;\mu,\Sigma)
    =
    \frac{1}{(2\pi)^{r/2}|\Sigma|^{1/2}}
    \exp\left\{
        -\frac{1}{2}(z-\mu)^\top\Sigma^{-1}(z-\mu)
    \right\},
    \label{eq:gaussian-density}
\end{equation}
with $\mu\in\mathbb{R}^r$ and positive-definite $\Sigma\in\mathbb{R}^{r\times r}$.

Because the regime path is unobserved, direct maximisation of the complete-data likelihood is not feasible. The model is estimated with the Baum-Welch algorithm, an expectation-maximisation procedure for HMMs \citep{baum1970maximization,bilmes1998em}. In the E-step at iteration $q$, the smoothed state probabilities are computed:
\begin{equation}
    \gamma_t(k)
    =
    \Pr(R_t=k \mid Z_{1:T};\Theta^{(q)}),
    \label{eq:hmm-gamma}
\end{equation}
and the smoothed transition probabilities are
\begin{equation}
    \xi_t(j,k)
    =
    \Pr(R_{t-1}=j,R_t=k \mid Z_{1:T};\Theta^{(q)}),
    \qquad
    t=2,\ldots,T.
    \label{eq:hmm-xi}
\end{equation}
In the M-step, the parameters are updated by weighted averages:
\begin{align}
    \delta_k^{(q+1)}
    &=
    \gamma_1(k),
    \label{eq:update-delta}
    \\
    p_{jk}^{(q+1)}
    &=
    \frac{\sum_{t=2}^{T}\xi_t(j,k)}
         {\sum_{t=2}^{T}\sum_{\ell=1}^{2}\xi_t(j,\ell)},
    \label{eq:update-transition}
    \\
    \mu_k^{(q+1)}
    &=
    \frac{\sum_{t=1}^{T}\gamma_t(k)Z_t}
         {\sum_{t=1}^{T}\gamma_t(k)},
    \label{eq:update-mean}
    \\
    \Sigma_k^{(q+1)}
    &=
    \frac{\sum_{t=1}^{T}\gamma_t(k)
    (Z_t-\mu_k^{(q+1)})(Z_t-\mu_k^{(q+1)})^\top}
    {\sum_{t=1}^{T}\gamma_t(k)}.
    \label{eq:update-covariance}
\end{align}
The E-step and M-step are iterated until the observed-data log-likelihood stabilises:
\begin{equation}
    \ell(\Theta)
    =
    \log
    \left[
        \sum_{R_1=1}^{2}
        \cdots
        \sum_{R_T=1}^{2}
        f(Z_{1:T},R_{1:T};\Theta)
    \right].
    \label{eq:hmm-observed-loglikelihood}
\end{equation}

For downstream PD prediction, the relevant macro signal is not the smoothed probability in Equation~\eqref{eq:hmm-gamma}. Smoothed probabilities condition on the full factor path $Z_{1:T}$ and therefore use future observations. The PD model instead uses the forward-filtered probability conditional on the fitted full-window parameters:
\begin{equation}
    \alpha_t(k)
    =
    f(Z_{1:t},R_t=k;\hat{\Theta}),
    \label{eq:forward-variable}
\end{equation}
with recursion
\begin{align}
    \alpha_1(k)
    &=
    \hat{\delta}_k
    \phi_r(Z_1;\hat{\mu}_k,\hat{\Sigma}_k),
    \label{eq:forward-initialisation}
    \\
    \alpha_t(k)
    &=
    \phi_r(Z_t;\hat{\mu}_k,\hat{\Sigma}_k)
    \sum_{j=1}^{2}\alpha_{t-1}(j)\hat{p}_{jk},
    \qquad
    t=2,\ldots,T.
    \label{eq:forward-recursion}
\end{align}
\begin{equation}
    \hat p_t(k)
    =
    \Pr(R_t=k \mid Z_{1:t};\hat{\Theta})
    =
    \frac{\alpha_t(k)}
    {\sum_{\ell=1}^{2}\alpha_t(\ell)}.
    \label{eq:filtered-probability}
\end{equation}
The stress input used in the credit model is
\begin{equation}
    \hat p_t^{\mathrm{Stress}}
    =
    \hat p_t(k^\star),
    \qquad
    k^\star
    =
    \arg\max_{k\in\{1,2\}} s(k),
    \label{eq:stress-state-selection}
\end{equation}
where $s(k)$ is an economic stress score assigned from the fitted factor profile. Since every factor is oriented so that larger values indicate greater stress pressure, the study uses
\begin{equation}
    s(k)
    =
    \sum_{g=1}^{r}\mu_{k,g},
    \qquad
    k=1,2.
    \label{eq:stress-score}
\end{equation}
Thus, the stress state is the fitted state with the larger combined mean across the adverse-oriented factors.

\section{Model Comparison}

The thesis centres on three logistic specifications: a borrower-only baseline, a macro-augmented baseline, and a clean regime-aware model. The implemented nested ladder adds three diagnostic specifications so the incremental value of the stress level, stress interactions, raw macro variables, and stress dynamics can be separated while preserving a common sample and target.

\begin{table}[htbp]
    \centering
    \caption{Model hierarchy used in the empirical comparison}
    \label{tab:model-hierarchy}
    \begin{tabular}{L{0.22\textwidth}L{0.34\textwidth}L{0.29\textwidth}}
        \toprule
        Model & Specification & Purpose \\
        \midrule
        Borrower-only logit & Borrower and loan features $X_{i,t}$ only. & Establish the static loan-level benchmark. \\
        Macro-augmented logit & Borrower features plus transformed raw macro variables $M_t$. & Test whether observed macro variables improve PD estimation additively. \\
        Baseline plus stress & Borrower features plus $q_t$ without interactions. & Isolate the direct level information in the latent state. \\
        Regime-aware logit & Borrower features, filtered stress probability, and stress-by-borrower interactions. & Test whether latent stress changes borrower risk sensitivities. \\
        Regime-aware plus raw macro & Clean regime-aware specification plus $M_t$. & Test whether the HMM signal adds information beyond observed macro covariates. \\
        Dynamic-stress robustness & Clean regime-aware specification plus $\Delta q_t$ and stress duration. & Test a small dynamic extension without redefining the core model. \\
        \bottomrule
    \end{tabular}
\end{table}

The borrower-only baseline is
\begin{equation}
    \pi_{i,t}^{(0)}
    =
    \Pr(Y_{i,t}^{(12)}=1\mid X_{i,t})
    =
    \sigma\left(\beta_0 + X_{i,t}^{\top}\beta\right),
    \label{eq:borrower-baseline}
\end{equation}
where $\sigma(a)=(1+\exp(-a))^{-1}$. This model assumes that borrower and loan characteristics have fixed effects through time.

The macro-augmented baseline is
\begin{equation}
    \pi_{i,t}^{(M)}
    =
    \sigma\left(
        \beta_0
        + X_{i,t}^{\top}\beta
        + M_t^{\top}\gamma_M
    \right).
    \label{eq:macro-baseline}
\end{equation}
This model allows the macroeconomic environment to shift the PD, but it does not allow the effect of borrower variables to change with the macro state.

The proposed regime-aware model is
\begin{equation}
    \pi_{i,t}^{(R)}
    =
    \sigma\left(
        \beta_0
        + X_{i,t}^{\top}\beta
        + \gamma_R \hat p_t^{\mathrm{Stress}}
        + \hat p_t^{\mathrm{Stress}} X_{i,t}^{\top}\theta
    \right).
    \label{eq:regime-aware-model}
\end{equation}
The coefficient vector $\beta$ measures baseline borrower effects. The scalar $\gamma_R$ measures the direct level effect of macro stress probability. The interaction vector $\theta$ measures how borrower sensitivities change as the economy moves toward the stress regime.

To see the interpretation of the interaction terms, consider a continuous borrower feature $X_{i,t,j}$. The marginal effect on the log-odds is
\begin{equation}
    \frac{\partial}{\partial X_{i,t,j}}
    \operatorname{logit}\left(\pi_{i,t}^{(R)}\right)
    =
    \beta_j
    +
    \hat p_t^{\mathrm{Stress}}\theta_j.
    \label{eq:regime-marginal-effect}
\end{equation}
When $\theta_j>0$, the log-odds contribution of feature $j$ becomes stronger in stress. When $\theta_j<0$, the contribution weakens as stress probability rises. For binary or categorical indicators, the same expression is interpreted as a stress-dependent discrete log-odds difference rather than a derivative. This is the main mathematical distinction between the regime-aware model and the macro-augmented baseline.

The dynamic robustness variant adds
\begin{equation}
    \Delta q_t=q_t-q_{t-1},
    \qquad
    D_t^{\mathrm{Stress}}=
    \begin{cases}
       D_{t-1}^{\mathrm{Stress}}+1,&q_t\geq 0.5,\\
       0,&q_t<0.5,
    \end{cases}
    \label{eq:stress-dynamics}
\end{equation}
to the clean regime-aware model. Lagged probabilities at 3, 6, and 12 months are generated for diagnostics, but the reported smaller robustness model uses only $\Delta q_t$ and $D_t^{\mathrm{Stress}}$.

The logistic coefficients are estimated by penalised maximum likelihood. For a generic model with predicted probabilities $\hat{\pi}_{i,t}$, the Bernoulli log-likelihood over the training set $\mathcal{T}$ is
\begin{equation}
    \ell
    =
    \sum_{(i,t)\in\mathcal{T}}
    \left[
        Y_{i,t}^{(12)}\log(\hat{\pi}_{i,t})
        +
        \left(1-Y_{i,t}^{(12)}\right)
        \log(1-\hat{\pi}_{i,t})
    \right].
    \label{eq:logit-likelihood}
\end{equation}
All numerical features and constructed interactions are standardised inside the pipeline. The empirical implementation uses stochastic-gradient logistic regression with an $L_2$ penalty on the complete non-intercept coefficient vector:
\begin{equation}
    \widehat{\Omega}
    =
    \arg\min_{\Omega}
    \left\{
        -\ell(\Omega)
        +
        \lambda \|\omega\|_2^2
    \right\},
    \qquad
    \lambda\geq 0,
    \label{eq:ridge-objective}
\end{equation}
where $\Omega=(\beta_0,\omega)$ and $\omega$ contains every non-intercept coefficient in the relevant specification. The penalty stabilises correlated macro and interaction features \citep{lecessie1992ridge}. The implementation evaluates $\lambda\in\{3\times10^{-6},10^{-5},3\times10^{-5}\}$ on deterministic 500,000-row training and validation subsamples, selects the value with the lowest validation Brier score, and then fits the selected pipeline on the full training sample.

\section{Evaluation Plan}

All comparisons use chronological splits. This is necessary because the task is forward-looking: a credit model trained with future information would overstate performance. Let $\tau_1<\tau_2$ denote calendar-month cutoffs. The full modelling sample is partitioned into training, validation, and test periods as
\begin{equation}
    \begin{aligned}
    \mathcal{R}_{\mathrm{train}}
    &=
    \{(i,t)\in\mathcal{R}: t\leq\tau_1\},\\
    \mathcal{R}_{\mathrm{valid}}
    &=
    \{(i,t)\in\mathcal{R}: \tau_1<t\leq\tau_2\},\\
    \mathcal{R}_{\mathrm{test}}
    &=
    \{(i,t)\in\mathcal{R}: t>\tau_2\}.
    \end{aligned}
    \label{eq:chronological-split}
\end{equation}
The current empirical pass estimates one HMM over the mortgage study window plus a 12-month lookback. The forward recursion ensures that the state update at $t$ does not use factor observations after $t$, but its fitted scaler and HMM parameters use the full study window. This is less forward-looking than smoothing the state sequence, but it is not a strict real-time design. The required publication robustness pass will re-estimate the factor transforms and HMM on expanding historical windows. The downstream PD models are fit on $\mathcal{R}_{\mathrm{train}}$, tuned on $\mathcal{R}_{\mathrm{valid}}$, calibrated on validation predictions, and evaluated once on $\mathcal{R}_{\mathrm{test}}$ in the present results.

The primary evaluation criterion is calibration. The Brier score is
\begin{equation}
    BS
    =
    \frac{1}{|\mathcal{S}|}
    \sum_{(i,t)\in\mathcal{S}}
    \left(\hat{\pi}_{i,t}-Y_{i,t}^{(12)}\right)^2,
    \label{eq:brier-score}
\end{equation}
where $\mathcal{S}$ denotes the evaluation sample and $|\mathcal{S}|$ is its number of loan-month observations. Lower values indicate better probabilistic accuracy \citep{brier1950}.

The Brier score can be decomposed into reliability, resolution, and uncertainty. For binned calibration diagnostics, let predictions be grouped into $B$ bins, with $n_b$ observations in bin $b$, $n=|\mathcal{S}|$, average predicted probability $\bar{\pi}_b$, observed default rate $\bar{y}_b$, and overall default rate $\bar{y}$. The grouped Murphy decomposition is
\begin{equation}
    BS_{\mathrm{grouped}}
    =
    \underbrace{\sum_{b=1}^{B}\frac{n_b}{n}(\bar{\pi}_b-\bar{y}_b)^2}_{\mathrm{Reliability}}
    -
    \underbrace{\sum_{b=1}^{B}\frac{n_b}{n}(\bar{y}_b-\bar{y})^2}_{\mathrm{Resolution}}
    +
    \underbrace{\bar{y}(1-\bar{y})}_{\mathrm{Uncertainty}}.
    \label{eq:brier-decomposition}
\end{equation}
Reliability measures calibration error. Resolution measures the ability of the model to separate groups with different observed default frequencies. Uncertainty depends only on the base default rate \citep{murphy1973vector,dimitriadis2020_reliability}. The grouped decomposition is a diagnostic summary; the ungrouped Brier score in Equation~\ref{eq:brier-score} remains the primary reported score.

Calibration is also evaluated with bin-level observed-versus-predicted default rates:
\begin{equation}
    \widehat{\mathrm{Cal}}_b
    =
    \left(
        \frac{1}{n_b}\sum_{(i,t)\in b}\hat{\pi}_{i,t},
        \frac{1}{n_b}\sum_{(i,t)\in b}Y_{i,t}^{(12)}
    \right),
    \qquad
    b=1,\ldots,B.
    \label{eq:calibration-points}
\end{equation}
A well-calibrated model should produce bin-level points with small vertical deviations from the 45-degree line; the reliability term in Equation~\ref{eq:brier-decomposition} summarises these deviations numerically. Calibration is examined overall, by calendar period, and by regime slice. The regime-slice check is central to the thesis because the proposed model is intended to improve probability quality during macro stress and transition periods.

Three validation-fitted calibration maps are compared. Platt scaling fits
\begin{equation}
    \widehat\pi^{\mathrm{Platt}}
    =\sigma\!\left(a+b\,\operatorname{logit}(\widehat\pi^{\mathrm{raw}})\right),
    \label{eq:platt-calibration}
\end{equation}
intercept-only recalibration fixes $b=1$, and isotonic regression estimates a non-decreasing map from raw to calibrated probabilities. The primary output uses Platt scaling for continuity across specifications; the other maps are reported as sensitivity analysis. Calibration parameters are estimated only from the chronological validation predictions.

Portfolio-level timing is assessed by the monthly predicted-default count
\begin{equation}
    \widehat D_t=\sum_{i\in\mathcal{R}_t}\widehat\pi_{i,t},
    \qquad
    D_t=\sum_{i\in\mathcal{R}_t}Y_{i,t}^{(12)},
    \label{eq:monthly-default-count}
\end{equation}
with absolute error $|\widehat D_t-D_t|$ and absolute relative error $|\widehat D_t-D_t|/D_t$. This diagnostic tests whether improvements in average loan-level scores translate into useful portfolio risk levels through time.

Secondary metrics include log loss, ROC-AUC, and PR-AUC. Log loss is
\begin{equation}
    LL
    =
    -
    \frac{1}{|\mathcal{S}|}
    \sum_{(i,t)\in\mathcal{S}}
    \left[
        Y_{i,t}^{(12)}\log(\hat{\pi}_{i,t})
        +
        \left(1-Y_{i,t}^{(12)}\right)\log(1-\hat{\pi}_{i,t})
    \right].
    \label{eq:log-loss}
\end{equation}
For numerical stability, predicted probabilities are clipped to $[\varepsilon,1-\varepsilon]$ before evaluating the logarithms, with $\varepsilon=10^{-15}$ fixed before testing.
ROC-AUC is retained as a ranking measure, while PR-AUC is reported because serious delinquency is rare and precision-recall analysis is informative for imbalanced binary outcomes \citep{davis_goadrich_2006,saito2015prroc}.

The empirical interpretation is based on the full pattern of evidence rather than one metric alone. The regime-aware model is considered useful only if probability-score gains are accompanied by credible temporal, calibration, and portfolio diagnostics. Any deterioration relative to the baselines is reported explicitly. This criterion follows the purpose of the model: to produce interpretable and macro-sensitive PD estimates, not merely to maximise a black-box accuracy measure.

\begin{table}[htbp]
    \centering
    \caption{Evaluation metrics and interpretation}
    \label{tab:evaluation-metrics}
    \begin{tabular}{L{0.22\textwidth}L{0.31\textwidth}L{0.31\textwidth}}
        \toprule
        Metric & Main question answered & Role in this thesis \\
        \midrule
        Brier score & What is the mean squared error of the predicted PDs? & Primary probability-quality metric. \\
        Reliability diagram & Do predicted PD bands match observed default rates? & Main visual calibration diagnostic. \\
        Brier decomposition & Are changes driven by reliability, resolution, or base-rate uncertainty? & Explains the source of model improvement. \\
        Log loss & Does the model assign high probability to realised outcomes? & Secondary proper scoring rule. \\
        ROC-AUC & Does the model rank defaulters above non-defaulters? & Secondary discrimination metric. \\
        PR-AUC & How well does the model identify rare default events? & Secondary rare-event metric. \\
        Monthly count error & Does the sum of predicted PDs track realised defaults by month? & Portfolio-level calibration guardrail. \\
        \bottomrule
    \end{tabular}
\end{table}
